\chapterrule
\section{线性联想记忆：写入、串扰与误差修正}
\label{intro:key-value-outer-product}
\label{intro:delta-memory}

\phantomsection\label{concept:outer-product}

\noindent\textbf{{\symboliconfont ★} 扩展路径：从存入与取出的角度重看 Attention。}
从主线已经学过的点积、Attention 与 FFN 出发，
本部分把同样的计算拆成外积写入、串扰、误差修正与非线性检索逐步复算。
若读过\hyperref[sec:modern-delta-rule]{现代架构章的选读 Delta 规则}，也可以在这里核对它的写入过程。
这里改写的是外部状态 $\mathcal{M}$；它与训练神经网络权重是两种更新。

Transformer 主线里每个 Token 占一行。现在换一种视角理解线性层：
能否把“遇见 Key 就读出 Value”的关联写入一个矩阵？
和前面一样，Key、Query 与 Value 都以\textbf{一条记录占一行}来表示。

\subsection{先问输入、输出和形状}
单次 Query 是 $\vect{x}^\top:1\times d_k$，Key 为 $\vect{k}^\top:1\times d_k$，
Value 为 $\vect{v}^\top:1\times d_v$。我们希望读出一行 $1\times d_v$ 的结果。
\[
\underbrace{\vect{x}^\top}_{1\times d_k}
\underbrace{\vect{k}}_{d_k\times1}
\underbrace{\vect{v}^\top}_{1\times d_v}
=(\vect{x}^\top\vect{k})\vect{v}^\top.
\]
先用点积得到匹配分数，再乘 Value；不能直接把 Key 和 Value 的维度混同。

\subsection{亲手写入一条二维记忆}
先写 A：令 $\vect{k}_A^\top=(1,0)$，$\vect{v}_A^\top=(1,0)$，查询 $\vect{x}^\top=\vect{k}_A^\top$。
下面先记作 $\vect{k}^\top=\vect{k}_A^\top,\vect{v}^\top=\vect{v}_A^\top$；后续加入 B 时继续使用这条 A。
两条向量的外积矩阵是
\[
\mathcal{M}=\vect{k} \vect{v}^\top
=\begin{bmatrix}1\\0\end{bmatrix}
\begin{bmatrix}1&0\end{bmatrix}
=\boxed{\begin{bmatrix}1&0\\0&0\end{bmatrix}}.
\]
读取时从左边乘一条\textbf{行向量}：
\begin{equation}
\boxed{\vect{x}^\top\mathcal{M}=(1,0)\mathcal{M}=(1,0).}
\end{equation}
因为 $\vect{k}^\top\vect{k}=1$，所以 $\vect{k}^\top\mathcal{M}=(\vect{k}^\top\vect{k})\vect{v}^\top=\vect{v}^\top$；
若 Key 并非单位向量，结果一般会乘上 $\|\vect{k}\|^2$。
外积得到一个矩阵，点积得到一个数；二者不要混淆。

\begin{idea}
这就是同一计算的两种写法：$\vect{x}^\top(\vect{k} \vect{v}^\top)=(\vect{x}^\top\vect{k})\vect{v}^\top$。
如果有 $T_k$ 条记忆，Key 和 Value 分别\textbf{按行}存进
$K:T_k\times d_k$、$V:T_k\times d_v$，
匹配权重为 $\vect{a}^\top:1\times T_k$，
则最后读出为 $\vect{a}^\top V:1\times d_v$。
形式直接对齐主线 Attention 的“权重乘 Value”。
\end{idea}

\subsection{多条记忆：为什么会互相干扰？}
\phantomsection\label{intro:memory-interference}
\phantomsection\label{concept:memory-interference}
\phantomsection\label{partq:memory}
对 $T_k$ 条行向量 Key--Value 配对，最直接的写入是外积累加：
\begin{equation}
\boxed{\mathcal{M}=\sum_{i=1}^{T_k}\vect{k}_i \vect{v}_i^\top},
\qquad
\boxed{\vect{x}^\top\mathcal{M}=\sum_{i=1}^{T_k}(\vect{x}^\top\vect{k}_i)\vect{v}_i^\top}.
\end{equation}
这里 $\mathcal{M}:d_k\times d_v$，$\vect{x}^\top:1\times d_k$，读取输出仍是一行。

\subsubsection{两条记忆的二维例子}
\[
\vect{k}_A^\top=(1,0),\quad \vect{v}_A^\top=(1,0),\qquad
\vect{k}_B^\top=(0.6,0.8),\quad \vect{v}_B^\top=(0,1).
\]
写入矩阵分别是
\[
\vect{k}_A \vect{v}_A^\top=\begin{bmatrix}1&0\\0&0\end{bmatrix},
\qquad
\vect{k}_B \vect{v}_B^\top=\begin{bmatrix}0&0.6\\0&0.8\end{bmatrix}.
\]
因此
\[
\boxed{\mathcal{M}_{\rm simple}=\begin{bmatrix}1&0.6\\0&0.8\end{bmatrix}}.
\]
查询 A：
\[
\vect{k}_A^\top \mathcal{M}_{\rm simple}
=(1,0)\begin{bmatrix}1&0.6\\0&0.8\end{bmatrix}
=\boxed{(1,0.6)}
=\underbrace{(1,0)}_{\text{正确信号}}
+\underbrace{(0,0.6)}_{\text{来自 B 的串扰}}.
\]
B 被激活，是因为 $\vect{k}_A^\top\vect{k}_B=0.6$ 不为零；
两条 Key 虽然都是单位长度，却不垂直。

对单位行向量 Key $\vect{k}_j^\top$（$\vect{k}_j^\top\vect{k}_j=1$），恒有
\[
\vect{k}_j^\top\mathcal{M}=\vect{v}_j^\top+\sum_{i\ne j}(\vect{k}_j^\top\vect{k}_i)\vect{v}_i^\top.
\]
其他 Key 的匹配分数不为零时，就会产生干扰。

\subsection{联立求解：同时拟合两条记忆}
直接累加外积时，每条新记忆可能改变旧 Key 的输出。
可以让矩阵同时拟合全部输入与目标：
\[
\boxed{\min_{\mathcal{M}}\sum_{i=1}^{T_k}\|\vect{k}_i^\top\mathcal{M}-\vect{v}_i^\top\|^2}.
\]
这叫最小二乘：让同一个状态同时满足各条关联。

刚才的外积相加得到
\[
\mathcal{M}_{\rm simple}=\begin{bmatrix}1&0.6\\0&0.8\end{bmatrix}.
\]
要同时读对 A、B，可以直接联立求解。设矩阵的两行为 $\vect{m}_1^\top,\vect{m}_2^\top$。
A 要求 $\vect{m}_1^\top=(1,0)$；B 要求 $0.6\vect{m}_1^\top+0.8\vect{m}_2^\top=(0,1)$，所以
\[
\vect{m}_2^\top=\frac{(0,1)-0.6(1,0)}{0.8}=(-0.75,1.25).
\]
于是
\[
\boxed{\mathcal{M}_{\rm better}=\begin{bmatrix}1&0\\-0.75&1.25\end{bmatrix}},
\]
就能逐项验算：
\[
(1,0)\mathcal{M}_{\rm better}=(1,0),\qquad
(0.6,0.8)\mathcal{M}_{\rm better}=(0,1).
\]
负权重在这里用于\textbf{抵消交叉干扰}。
但如果约束不相容，线性层并不能任意精确存取无限多条关联。

\subsection{从误差修正到 Delta 写入}
\phantomsection\label{concept:delta-rule}
前面已知道怎样让一个矩阵拟合目标。现在每次只来一条配对，
能否像 Mini-batch 训练一样，先看当前误差，再走一步负梯度？
对当前 Key--Value 配对，定义平方误差与读取残差：
\[
\ell_t(\mathcal{M})=\tfrac12\|\vect k_t^\top\mathcal{M}-\vect v_t^\top\|^2,
\qquad
\vect r_t^\top=\vect v_t^\top-\vect k_t^\top\mathcal{M}_{t-1}.
\]
残差就是“期望读出减去已有读出”。沿用
\hyperref[sec:mlp-matrix-gradient]{矩阵反向传播}的“输入列乘输出误差行”，有
\[
\frac{\partial\ell_t}{\partial\mathcal{M}}
=\vect k_t(\vect k_t^\top\mathcal{M}_{t-1}-\vect v_t^\top)
=-\vect k_t\vect r_t^\top.
\]
减去学习率乘梯度，便得到\textbf{Delta 写入}：
\[
\boxed{\mathcal{M}_t=\mathcal{M}_{t-1}+\eta\vect k_t\vect r_t^\top}.
\]
这里 $\vect k_t^\top:1\times d_k$、$\vect r_t^\top:1\times d_v$，
所以外积与状态同为 $d_k\times d_v$。
若 Key 长度为 $1$ 且 $\eta=1$，当次写入可消除该 Key 的读取误差，但也可能改变其他 Key 的读出。
形式上这是一次单样本梯度更新；此处更新对象是随记录变化的\textbf{记忆状态} $\mathcal{M}$，
与训练模型的投影权重分开。本例反复调整状态时，Key、Value 始终固定。

\subsection{［复算］写好 B 时，A 会不会再次出错？}

仍用两条单位 Key 和 Value：
\[
\vect{k}_A^\top=(1,0),\ \vect{v}_A^\top=(1,0),\qquad
\vect{k}_B^\top=(0.6,0.8),\ \vect{v}_B^\top=(0,1).
\]
从 $\mathcal{M}_0=0$、$\eta=1$ 开始，先写 A：
\[
\mathcal{M}_1=\vect{k}_A \vect{v}_A^\top=\begin{bmatrix}1&0\\0&0\end{bmatrix}.
\]
写 B 前，旧状态的读取为 $\vect{k}_B^\top\mathcal{M}_1=(0.6,0)$，
所以不是直接追加 $\vect{k}_B \vect{v}_B^\top$，而是计算残差：
\[
\vect{r}_B^\top=\vect{v}_B^\top-\vect{k}_B^\top\mathcal{M}_1=(-0.6,1).
\]
下一步得到
\begin{equation}
\begin{aligned}
\mathcal{M}_2&=\mathcal{M}_1+\vect{k}_B \vect{r}_B^\top\\
&=\begin{bmatrix}1&0\\0&0\end{bmatrix}
+\begin{bmatrix}0.6\\0.8\end{bmatrix}(-0.6,1)
=\boxed{\begin{bmatrix}0.64&0.6\\-0.48&0.8\end{bmatrix}}.
\end{aligned}
\label{eq:delta-overwrites-key-a}
\end{equation}
立即查询 B 得到\textbf{正确的} $\vect{k}_B^\top\mathcal{M}_2=(0,1)$；
反过来，旧的 A 查询却变成
$\vect{k}_A^\top\mathcal{M}_2=(0.64,0.6)$，不再等于 $(1,0)$。
A 的误差范数为
\[
\|\vect{k}_A^\top\mathcal{M}_2-\vect{v}_A^\top\|_2=\sqrt{(-0.36)^2+0.6^2}
\approx\boxed{0.700}.
\]
两种方法都从同一个“只写好 A”的 $\mathcal{M}_1$ 出发。
若采用直接相加外积来写 B，A 的误差为 $0.600$。
\textbf{在这个明确例子里，单步 Delta 对旧记忆的扰动反而更大。}

\subsection{把 A、B 再访问几轮，误差怎样变化？}
\label{sec:delta-revisit}

写好 B 后，A 又错了；下一次重新访问 A 时，就按 A 的新残差修正。
从上面的 $\mathcal{M}_2$ 出发，A 的残差为 $(0.36,-0.6)$，所以
\[
\mathcal{M}_3
=\begin{bmatrix}0.64&0.6\\-0.48&0.8\end{bmatrix}
+\begin{bmatrix}1\\0\end{bmatrix}(0.36,-0.6)
=\begin{bmatrix}1&0\\-0.48&0.8\end{bmatrix}.
\]
A 当次恢复正确，但 B 读成 $(0.216,0.64)$，还需按 B 的残差继续修正。
把从零开始依次访问一次 A、一次 B 叫作一轮；始终保持 $\eta=1$，结果如下。
\begin{center}\small
\tagpdfsetup{table/header-rows={1}}
\begin{tabular}{@{}rrr@{}}\toprule
完整 A$\to$B 轮数 & A 的误差范数 & B 的误差范数\\\midrule
1 & 0.699714 & 0\\
2 & 0.251897 & 0\\
3 & 0.090683 & 0\\
5 & 0.011753 & 0\\\bottomrule
\end{tabular}
\end{center}
每行都在写完 B 后测量，因此 B 恰好刚被修正，误差为零。
A 的误差随完整轮次缩小；在这两条线性无关的单位 Key、固定目标与 $\eta=1$ 的条件下，
反复重访收敛到前面联立求出的 $\mathcal{M}_{\rm better}$。
运行 \texttt{python3 scripts/memory\_demo.py} 可复算这张表。

\begin{boundarybox}
Delta 写入按残差修正当前读取；每次修正都可能改变旧记忆。
因此要在每次更新后检查全部 Key 的读出，并按完整轮次比较误差。
\end{boundarybox}

\begin{idea}
最小二乘是一个全局优化目标；Delta 写入是一种依赖当前状态的局部更新。
两者都在减少干扰：前者一次求联立解，后者逐条修正当前读出。
\hyperref[sec:modern-delta-rule]{现代架构章中的 Delta 写入}也按当前读出的残差更新状态；
Gated DeltaNet 则加入遗忘门控，调节旧记忆保留多少。
\end{idea}

\noindent\emph{对应迁移练习：}\hyperref[exercise:delta-reverse]{交换写入次序的迁移题}
交换 A、B 的写入次序，检查被扰动的是否始终是旧 Key。

\chapterchecktasks{key-value-outer-product}{理解、手算与判断}{
\label{check:delta-memory}
\item\textbf{形状：}行向量 Key $\vect{k}^\top:1\times2$、Value $\vect{v}^\top:1\times3$。外积写入 $\mathcal{M}=\vect{k} \vect{v}^\top$ 的形状是什么？查询 $\vect{q}^\top:1\times2$ 后 $\vect{q}^\top\mathcal{M}$ 又是什么形状？\chapterinlinehint{先写转置 Key 的形状，再核对内侧维度。}
\item\textbf{预测：}同一条记忆若把 Key 写入时改为 $\vect{k}'^\top=(0.5,0)$，再以 $\vect{q}^\top=\vect{k}'^\top$ 读取，恢复的 Value 是多少倍？\chapterinlinehint{把写入和读取两次出现的 Key 缩放都算上。}
\item\textbf{手算：}$\vect{k}_A^\top=(1,0),\vect{k}_B^\top=(0.6,0.8)$；$\vect{v}_A^\top=(1,0),\vect{v}_B^\top=(0,1)$。把两条外积相加，用 $\vect{k}_A^\top$ 查询得到什么？\chapterinlinehint{分别求两条外积，再相加并用查询读取。}
\item\textbf{预测：}若把 $\vect{k}_B^\top$ 改为 $(0,1)$，保持其他项不变，$\vect{k}_A^\top$ 查询时串扰如何变化？\chapterinlinehint{先比较修改前后的交叉点积。}
\item\textbf{定位：}Delta 写入为什么先用当前状态读取一次？“期望 Value 减去已有读出”代表什么量？\chapterinlinehint{先写目标读出与当前读出的差。}
\item\textbf{手算：}写入 B 后 $\mathcal{M}_2=\begin{bmatrix}0.64&0.6\\-0.48&0.8\end{bmatrix}$；用 $\vect{k}_A^\top=(1,0)$ 读出什么？与目标 $\vect{v}_A^\top=(1,0)$ 相差多少？\chapterinlinehint{先用查询选出对应行，再减去目标向量。}
\item\textbf{找错：}同学说“B 的残差为零后，原先的 A 也必定被完全修复”。指出更新为何可能影响其他查询。\chapterinlinehint{考察新增外积与另一条 Key 的点积。}
}

\par\smallskip
减少串扰的线性办法都在\textbf{写入端}：联立求解或逐条按残差修正。
下一章转到\textbf{读取端}，让匹配分数先经过非线性筛选再读出 Value。
